\section{异常控制流}

\subsection{异常}

异常 (exception) 就是控制流中的突变，用来响应处理器状态中的某些变化。

\begin{figure}[H]
    \centering
    \includegraphics[width=0.6\textwidth]{ECF/exception.png}
\end{figure}

在处理器中，状态被编码为不同的位和信号。状态变化称为事件 (event) 。 
在任何情况下，当处理器检测到有事件发生时，它就会通过一张叫做异常表 (exception table) 的跳转表，
进行一个间接过程调用（异常），到一个专门设计用来处理这类事件
的操作系统子程序（异常处理程序 (exception handler)）。当异常处理程序完成处理后，根
据引起异常的事件的类型，会发生以下 3 种情况中的一种：
\begin{itemize}
  \item 处理程序将控制返回给当前指令，即当事件发生时正在执行的指令。
  \item 处理程序将控制返回给下一条指令， 如果没有发生异常将会执行的下一条指令。
  \item 处理程序终止被中断的程序。
\end{itemize}

\subsubsection{异常处理}

系统中可能的每种类型的异常都分配了一个唯一的非负整数的异常号 (exception number) 。
其中一些号码是处理器的设计者分配的，其他号码是由操作系统内核（操作系统
常驻内存的部分）的设计者分配的。前者的示例包括被零除、缺页、内存访问违例、断点
以及算术运算溢出。后者的示例包括系统调用和来自外部 I/O 设备的信号。

\begin{figure}[H]
    \centering
    \includegraphics[width=0.6\textwidth]{ECF/exception-tables.png}
\end{figure}

在系统启动时（当计算机重启或者加电时），操作系统分配和初始化一张称为异常表
的跳转表，使得表目 k 包含异常 k 的处理程序的地址。
在运行时（当系统在执行某个程序时），处理器检测到发生了一个事件，并且确定了相应
的异常号 k 。随后，处理器触发异常，方法是执行间接过程调用，通过异常表的表目 k，转到相应的处理程序。

异常类似于过程调用，但是有一些重要的不同之处：
\begin{itemize}
    \item 过程调用时，在跳转到处理程序之前，处理器将返回地址压入栈中。然而，根据异
常的类型，返回地址要么是当前指令（当事件发生时正在执行的指令），要么是下一
条指令（如果事件不发生，将会在当前指令后执行的指令）。
    \item 处理器也把一些额外的处理器状态压到栈里，在处理程序返回时，重新开始执行被
中断的程序会需要这些状态。比如， x86-64 系统会将包含当前条件码的 EFLAGS
寄存器和其他内容压入栈中。
    \item 如果控制从用户程序转移到内核，所有这些项目都被压到内核栈中，而不是压到用户栈中。
    \item 异常处理程序运行在内核模式下，这意味着它们对所有的系统资源都有完全的访问权限。
\end{itemize}

一旦硬件触发了异常，剩下的工作就是由异常处理程序在软件中完成。在处理程序处
理完事件之后，它通过执行一条特殊的“从中断返回”指令，可选地返回到被中断的程
序，该指令将适当的状态弹回到处理器的控制和数据寄存器中，如果异常中断的是一个用
户程序，就将状态恢复为用户模式，然后将控制返回给被中断的程序。

\subsubsection{异常处理的类别}

异常可以分为四类：中断 (interrupt) 、陷阱 (trap) 、故障 (fault) 和终止 (abort) 。 

\begin{figure}[H]
    \centering
    \includegraphics[width=0.6\textwidth]{ECF/types-of-exception.png}
\end{figure}

\paragraph{中断}

中断是异步发生的，是来自处理器外部的 I/O 设备的信号的结果。硬件中断不是由任
何一条专门的指令造成的，从这个意义上来说它是异步的。硬件中断的异常处理程序常常
称为中断处理程序 (interrupt handler) 。

\begin{figure}[H]
    \centering
    \includegraphics[width=0.6\textwidth]{ECF/interrupt.png}
\end{figure}

\paragraph{陷阱和系统调用}

陷阱是有意的异常，是执行一条指令的结果。就像中断处理程序一样，陷阱处理程序
将控制返回到下一条指令。陷阱最重要的用途是在用户程序和内核之间提供一个像过程一
样的接口，叫做系统调用。

\begin{figure}[H]
    \centering
    \includegraphics[width=0.6\textwidth]{ECF/trap.png}
\end{figure}

普通的函数运行在用户模式中，用户模式限制了函数可以执行的指令的类型，而且它们
只能访问与调用函数相同的栈。系统调用运行在内核模式中，内核模式允许系统调用执行特
权指令，并访问定义在内核中的栈。

\paragraph{故障}

故障由错误情况引起，它可能能够被故障处理程序修正。当故障发生时，处理器将控
制转移给故障处理程序。如果处理程序能够修正这个错误情况，它就将控制返回到引起故
障的指令，从而重新执行它。否则，处理程序返回到内核中的 abort 例程， abort 例程会
终止引起故障的应用程序。

\begin{figure}[H]
    \centering
    \includegraphics[width=0.6\textwidth]{ECF/fault.png}
\end{figure}

\paragraph{终止}

终止是不可恢复的致命错误造成的结果，通常是一些硬件错误，比如 DRAM 或者
SRAM 位被损坏时发生的奇偶错误。终止处理程序从不将控制返回给应用程序。

\begin{figure}[H]
    \centering
    \includegraphics[width=0.6\textwidth]{ECF/abort.png}
\end{figure}

\subsection{进程}

异常是允许操作系统内核提供进程 (process) 概念的基本构造块，进程是计算机科学中最深刻、最成功的概念之一。

\subsubsection{逻辑控制流}

即使在系统中通常有许多其他程序在运行，进程也可以向每个程序提供一种假象，好
像它在独占地使用处理器。如果想用调试器单步执行程序，我们会看到一系列的程序计数
器 (PC) 的值。这个 PC 值的序列叫做逻辑控制流，或者简称逻辑流。

\subsubsection{并发流}

一个逻辑流的执行在时间上与另一个流重叠，称为并发流 (concurrent flow) ，这两个
流被称为并发地运行。更准确地说，流 X 和 Y 互相并发，当且仅当 X 在 Y 开始之后和 Y
结束之前开始，或者 Y 在 X 开始之后和 X 结束之前开始。 

\begin{figure}[H]
    \centering
    \includegraphics[width=0.6\textwidth]{ECF/concurrency.png}
\end{figure}

多个流并发地执行的一般现象被称为并发 (concurrency) 。 一个进程和其他进程轮流运
行的概念称为多任务 (multitasking) 。 一个进程执行它的控制流的一部分的每一时间段叫
做时间片 (tim e slice) 。因此，多任务也叫做时间分片 (time slicing) 。

如果两个流 并发地运行在不同的处理器核或者计算机上，那么我们称它们为并行流 (parallel flow)。

\subsubsection{私有地址空间}

进程也为每个程序提供一种假象，好像它独占地使用系统地址空间。在一台 n 位地址
的机器上，地址空间是 $2^n$ 个可能地址的集合，0，1，...，$2^n - 1$ 。 进程为每个程序提供它
自己的私有地址空间。一般而言 ，和这个空间中某个地址相关联的那个内存字节是不能被
其他进程读或者写的，从这个意义上说，这个地址空间是私有的。

尽管和每个私有地址空间相关联的内存的内容一般是不同的，但是每个这样的空间都有相同的通用结构。

\begin{figure}[H]
    \centering
    \includegraphics[width=0.6\textwidth]{ECF/address-space.png}
\end{figure}

地址空间底部是保留给用户程序的，包括通常的代码、数据、堆和栈段。代码段总是
从地址 0x400000 开始。地址空间顶部保留给内核（操作系统常驻内存的部分）。地址空间
的这个部分包含内核在代表进程执行指令时（比如当应用程序执行系统调用时）使用的代
码、数据和栈。

\subsubsection{用户模式和内核模式}

处理器通常是用某个控制寄存器中的一个模式位 (mode bit) 来提供实现模式的切换，该寄
存器描述了进程当前享有的特权。当设置了模式位时，进程就运行在内核模式中（有时叫
做超级用户模式）。一个运行在内核模式的进程可以执行指令集中的任何指令，并且可以
访问系统中的任何内存位置。

没有设置模式位时，进程就运行在用户模式中。用户模式中的进程不允许执行特权指令
(privileged instruction)， 比如停止处理器、改变模式位，或者发起一个 I/O 操作。也不允许
用户模式中的进程直接引用地址空间中内核区内的代码和数据。任何这样的尝试都会导致致
命的保护故障。反之，用户程序必须通过系统调用接口间接地访问内核代码和数据。

\subsubsection{上下文切换}

操作系统内核使用一种称为上下文切换 (context switch) 的较高层形式的异常控制流来实现多任务。

内核为每个进程维持 一个上下文 (context) 。上下文就是内核重新启动一个被抢占的进
程所需的状态。它由一些对象的值组成，这些对象包括通用目的寄存器、浮点寄存器、程
序计数器、用户栈、状态寄存器、内核栈和各种内核数据结构，比如描述地址空间的页
表、包含有关当前进程信息的进程表，以及包含进程已打开文件的信息的文件表。

在进程执行的某些时刻，内核可以决定抢占当前进程，并重新开始一个先前被抢占了
的进程。这种决策就叫做调度 (scheduling)， 是由内核中称为调度器 (scheduler) 的代码处
理的。当内核选择一个新的进程运行时，我们说内核调度了这个进程。在内核调度了一个
新的进程运行后，它就抢占当前进程，并使用一种称为上下文切换的机制来将控制转移到
新的进程：保存当前进程的上下文；恢复某个先前被抢占的进程被保存的
上下文；将控制传递给这个新恢复的进程。

\subsection{系统调用错误处理}

当 Unix 系统级函数遇到错误时，它们通常会返回 -1，并设置全局整数变量 errno 来表示什么出错了。

通过使用错误处理包装函数，我们可以更可能地简化代码。包装函数调用基本函数，检查错误，如果有任何问题就终止。

\begin{lstlisting}[style=CStyle]
pid_t Fork(void)
{
    pid_t pid;
    if ((pid = fork()) < 0)
        unix_error("Fork error");
    return pid;
}
\end{lstlisting}

\subsection{进程控制}

\subsubsection{获取进程 ID}

每个进程都有一个唯一的正的进程 ID (PID) 。 getpid 函数返回调用进程的 PID。
getppid 函数返回它的父进程的 PID。

\begin{lstlisting}[style=CStyle]
#include <sys/types.h>
#include <unistd.h>

pid_t getpid(void);
pid_t getppid(void);
\end{lstlisting}

\subsubsection{创建和终止进程}

从程序员的角度，我们可以认为进程总是处于下面三种状态之一：

\begin{itemize}
    \item 运行：进进程要么在 CPU 上执行，要么在等待被执行且最终会被内核调度。
    \item 停止：进程的执行被挂起 (suspended)， 且不会被调度。当收到 SIGSTOP 、 SIGTSTP 、 SIGTTIN 或者 SIGTTOU 信号时，进程就停止，并且保持停止直到它收到
一个 SIGCONT 信号，在这个时刻，进程再次开始运行。
    \item 终止：进程永远地停止了。进程会因为三种原因终止：收到一个信号，该信号的
默认行为是终止进程；从主程序返回；调用 exit 函数。
\end{itemize}

\begin{lstlisting}[style=CStyle]
#include <stdlib.h>

void exit(int status);
\end{lstlisting}

exit 函数以 status 退出状态来终止进程。

\begin{lstlisting}[style=CStyle]
#include <sys/types.h>
#include <unistd.h>

pid_t fork(void);
\end{lstlisting}

父进程通过调用 fork 函数创建一个新的运行的子进程。

新创建的子进程几乎但不完全与父进程相同。子进程得到与父进程用户级虚拟地址空间相
同的（但是独立的）一份副本，包括代码和数据段、堆、共享库以及用户栈。子进程还获得与父
进程任何打开文件描述符相同的副本，这就意味着当父进程调用 fork 时，子进程可以读写父
进程中打开的任何文件。父进程和新创建的子进程之间最大的区别在于它们有不同的 PID。

fork 函数只被调用一次，却会返回两次：一次
是在调用进程（父进程）中，一次是在新创建的子进程中。在父进程中， fork 返回子进程
的 PID 。在子进程中， fork 返回 0 。因为子进程的 PID 总是为非零，返回值就提供一个明
确的方法来分辨程序是在父进程还是在子进程中执行。

\subsubsection{回收子进程}

一个进程由于某种原因终止时，内核并不是立即把它从系统中清除。相反，进程被
保持在一种已终止的状态中，直到被它的父进程回收 (reaped) 。当父进程回收己终止的子
进程时，内核将子进程的退出状态传递给父进程，然后抛弃己终止的进程，从此时开始，
该进程就不存在了。一个终止了但还未被回收的进程称为僵死进程 (zombie) 。

如果一个父进程终止了，内核会安排 init 进程成为它的孤儿进程的养父。 init 进程
的 PID 为 1，是在系统启动时由内核创建的，它不会终止，是所有进程的祖先。如果父进
程没有回收它的僵死子进程就终止了，那么内核会安排 init 进程去回收它们。

一个进程可以通过调用 waitpid 函数来等待它的子进程终止或者停止。

\begin{lstlisting}[style=CStyle]
#include <sys/types.h>
#include <sys/wait.h>

pid_t waitpid(pid_t pid， int *statusp， int options);
\end{lstlisting}

\paragraph{判定等待集合的成员}

等待集合的成员是由参数 pid 来确定的：
\begin{itemize}
    \item 如果 pid 大于 0，那么等待集合就是一个特定的子进程，ID 为 pid。
    \item 如果 pid 等于 -1，那么等待集合就是由父进程所有的子进程组成的。
\end{itemize}

\paragraph{修改默认行为}

认情况下（当 options=0 时）， waitpid 挂起调用进程的
执行，直到它的等待集合 (wait set) 中的一个子进程终止。如果等待集合中的一个进程在
刚调用的时刻就已经终止了，那么 waitpid 就立即返回。在这两种情况中， waitpid 返
回导致 waitpid 返回的已终止子进程的 PID 。此时，已终止的子进程已经被回收，内核会
从系统中删除掉它的所有痕迹。

可以通过设置 options 来修改默认行为：
\begin{itemize}
    \item WNOHANG: 如果等待集合中的任何子进程都还没有终止，那么就立即返回（返回值为 0）。
    \item WUNTRACED: 挂起调用进程的执行，直到等待集合中的一个进程变成巳终止或者被停止。返回的 PID 为导致返回的巳终止或被停止子进程的 PID 。
    \item WCONTINUED: 挂起调用进程的执行，直到等待集合中一个正在运行的进程终止或等待集合中一个被停止的进程收到 SIGCONT 信号重新开始执行。
    \item WNOHANG | WUNTRACED: 立即返回，如果等待集合中的子进程都没有被停止或终止，则返回值为 0；如果有一个停止或终止，则返回值为该子进程的 PID 。 
\end{itemize}

\paragraph{检查己回收子进程的退出状态}

如果 statusp 参数是非空的，那么 waitpid 就会在 status 中放上关千导致返回的
子进程的状态信息， status 是 statusp 指向的值。
\begin{itemize}
    \item WIFEXITED(status): 如果子进程通过调用 exit 或者一个返回 (return) 正常终止，就返回真。
    \item WEXITSTATUS (status): 只有在 WIFEXITED() 返回为真时，才会定义这个状态。返回一个正常终止的子进程的退出状态。
    \item WIFSIGNALED(status): 如果子进程是因为一个未被捕获的信号终止的，那么就返回真。
    \item WTERMSIG(status): 只有在 WIFSIGNALED() 返回为真时，才定义这个状态。返回导致子进程终止的信号的编号。
    \item WIFSTOPPED(status): 如果引起返回的子进程当前是停止的，那么就返回真。
    \item WSTOPSIG(status): 只有在 WIFSTOPPED() 返回为真时，才定义这个状态。返回引起子进程停止的信号的编号。
    \item WIFCONTINUED(status): 如果子进程收到 SIGCONT 信号重新启动，则返回真。
\end{itemize}

\paragraph{错误条件}

如果调用进程没有子进程，那么 waitpid 返回 -1， 并且设置 errno 为 ECHILD 。如果 waitpid 函数被一个信号中断，那么它返回 -1， 并设置 errno 为 EINTR 。

\subsubsection{让进程休眠}

sleep 函数将一个进程挂起一段指定的时间。

\begin{lstlisting}[style=CStyle]
#include <unistd.h>

unsigned int sleep(unsigned int secs);
\end{lstlisting}

如果请求的时间量已经到了，sleep 返回 0，否则返回还剩下的要休眠的秒数。

pause 函数让调用函数休眠，直到该进程收到一个信号。

\begin{lstlisting}[style=CStyle]
#include <unistd.h>

int pause(void);
\end{lstlisting}

\subsubsection{加载并运行程序}

execve 函数在当前进程的上下文中加载并运行一个新程序。

\begin{lstlisting}[style=CStyle]
#include <unistd.h>

int execve(const char *filename, char *const argv[], char *const envp[]);
\end{lstlisting}

execve 函数加载并运行可执行目标文件 filename, 且带参数列表 argv 和环境变量
列表 envp 。只有当出现错误时，例如找不到 filename, execve 才会返回到调用程序。
所以，与 fork 一次调用返回两次不同， execve 调用一次并从不返回。

\begin{figure}[H]
    \centering
    \includegraphics[width=0.6\textwidth]{ECF/argv.png}
    \includegraphics[width=0.6\textwidth]{ECF/envp.png}
\end{figure}

在 execve 加载了 filename 之后 ，调用启动代码。启动代码设置栈，并将控制传递给新程序的主函数。

\begin{figure}[H]
    \centering
    \includegraphics[width=0.6\textwidth]{ECF/program-stack-structure.png}
\end{figure}

\begin{sidenote}{程序和进程}
程序是一堆代码和数据；程序可以作为目标文件存在于磁盘上，或者作为段存在于地址空间
中。进程是执行中程序的一个具体的实例；程序总是运行在某个进程的上下文中。
fork 函数在新的子进程中运行相同的程序，新的子进程是父进程的一个复制品。
execve 函数在当前进程的上下文中加载并运行一个新的程序。
它会覆盖当前进程的地址空间，但并没有创建一个新进程。 
\end{sidenote}
















































\newpage